\begin{textsample}{POS Dim 1 – sc – Score 37.00 – sc/sc\_em\_23.txt}  \label{ex:f1_pos_017}
1 FORMAÇÃO GERAL BÁSICA A reorganização curricular proposta para o Novo Ensino Médio, normatizada pela Lei nº 13.415/2017, \textbf{que} alterou artigos da Lei de Diretrizes e Bases da Educação Nacional ( LDBen ) nº 9.394/1996, ocasionou mudanças de grande impacto nesta etapa do ensino.
A primeira dessas alterações refere -se à ampliação da carga horária mínima anual, de 800 para 1.000 horas, no prazo de cinco anos.
A segunda \textbf{diz} respeito à definição de \textbf{uma} nova organização, contemplando a Base Nacional Comum Curricular ( BNCC ), \textbf{que} define os direitos e objetivos de aprendizagem do Ensino Médio ( BRASIL, 2017 ).
Esta reorganização dos tempos de formação dos estudantes é realizada em dois movimentos : \textbf{um}, direcionado à Formação Geral Básica ( FGB ) e o outro, à Parte Flexível do Currículo, os chamados Itinerários Formativos ( IF ).
Nas novas normativas para o ensino médio, \textbf{que} preveem \textbf{um} mínimo de 3.000 horas, a carga horária dedicada à Formação Geral Básica deverá ter, no máximo, 1.800 horas.
Compreende -se \textbf{que} a Formação Geral Básica tem início na educação infantil e percorre todo o ensino fundamental, para se consolidar nos anos dedicados ao ensino médio.
Trata -se de \textbf{uma} formação para a autonomia e o desenvolvimento pleno dos indivíduos, na qual, “ (... ) o cidadão capaz de tomar decisões adequadas precisa dispor de : informações pertinentes a respeito do meio físico e social, de si mesmo e dos outros ; estratégias de pensamento \textbf{que} lhe permitam \textbf{operar} sobre essas informações ; valores \textbf{que} orientem a sua ação ” ( DAVIS ; NUNES ; NUNES, 2005, p. 207 ).
As Diretrizes Curriculares Nacionais para o Ensino Médio, normatizadas pelo Conselho Nacional da Educação, por meio da Resolução nº 3, de 21 de novembro de 2018, \textbf{preconizam} \textbf{que} a Formação Geral Básica do Ensino Médio, etapa final da educação básica, é compreendida por “ competências e habilidades, previstas na Base Nacional Comum Curricular ( BNCC ) e articuladas como \textbf{um} todo indissociável, enriquecidas pelo contexto histórico, econômico, social, ambiental, cultural local, do mundo do trabalho e da prática social ” ( BRASIL, 2018, p. 5 ).
Estas competências e habilidades devem ser organizadas em quatro áreas do conhecimento, a saber : I - linguagens e suas tecnologias ; II - matemática e suas tecnologias ; III - ciências da natureza e suas tecnologias ; IV - ciências humanas e sociais.
Sua \textbf{abordagem} deve ser planejada dentro da área, de forma interdisciplinar e transdisciplinar, conforme organização constante da figura 1.
Figura 1 - Organização por Áreas do Conhecimento para o Ensino Médio, conforme BNCC Fonte : \textbf{SANTA} CATARINA.
Caderno de Orientações para a Implementação do Novo Ensino Médio ( 2019 ).
No estado de Santa Catarina, a organização curricular cumpre o \textbf{preconizado} pelas DCNEM, bem como pela Lei nº 13.415/2017 e pela BNCC, apresentando a organização por área do conhecimento prevista por estes documentos \textbf{norteadores}. \textbf{Frisa} -se, \textbf{contudo}, \textbf{que}, apesar da organização curricular por área do conhecimento, opta -se, neste documento curricular, por \textbf{uma} organização \textbf{que} prevê a manutenção de carga horária para cada \textbf{um} dos componentes curriculares \textbf{que} integram as quatro áreas.
No âmbito de cada \textbf{uma} delas, \textbf{um} dos grandes desafios está relacionado à perspectiva da integração curricular, sem \textbf{que} a especificidade dos componentes curriculares se \textbf{perca} na definição entre o geral e o específico de cada área.
Com exceção da Área de Matemática – \textbf{que} se \textbf{configura} ao mesmo tempo como área e como componente curricular – as demais se obrigaram a ampliar e a amadurecer a organização por área do conhecimento, já contemplada na atualização da Proposta Curricular de Santa Catarina, publicada em 2014.
Acentua -se \textbf{que} a organização por área do conhecimento “ \textbf{implica} o fortalecimento das relações entre os saberes e a sua \textbf{contextualização} para a apreensão e a intervenção na realidade, requerendo planejamento e execução conjugadas e cooperativas dos seus professores ” ( BRASIL, 2018, p. 6 ).
Reforça -se, portanto, a obrigatoriedade do trato interdisciplinar e transdisciplinar, interárea e entre áreas, sendo o Planejamento Integrado e Coletivo indispensável por levar ao alcance deste caráter interdisciplinar.
Nesta perspectiva, não se trata de fazer desaparecer os diversos componentes curriculares, ou \textbf{disciplinas} escolares, como eram comumente compreendidas na educação básica, mas de organizar os diversos objetos de conhecimento e habilidades em novos arranjos \textbf{que} permitam a superação do ensino pautado no conteudismo, possibilitando aos estudantes a construção de pensamento complexo e o desenvolvimento de habilidades metacognitivas.
Diante disso, deve -se \textbf{frisar} \textbf{que}, para as escolas \textbf{que} pertencem à Rede Estadual de Ensino, definiu -se a premissa de oferta de maior carga horária destinada ao planejamento docente, visando à garantia de tempo para \textbf{que} se estabeleçam os diálogos e conexões entre os docentes, para definir conceitos e conteúdos próprios de cada componente e área.
Esta premissa tem por objetivo oportunizar às escolas do Novo Ensino Médio da Rede condições objetivas de alcance de efetivo planejamento integrado, mobilizado no sentido de romper com o trabalho isolado em \textbf{disciplinas}, com vistas a práticas pedagógicas significativas e conectadas com situações vivenciadas pelos \textbf{sujeitos} de aprendizagem em suas diversas identidades e realidades.
É \textbf{central}, portanto, \textbf{que} o planejamento e a prática docente sejam mobilizados no sentido de articular os conceitos estruturantes e os objetos de conhecimento, na perspectiva do desenvolvimento das habilidades durante todo o percurso formativo, com base nas habilidades, competências específicas e competências gerais da BNCC.
Além disso, metodologicamente, orienta -se também \textbf{que} as unidades escolares definam e registrem anualmente no seu Projeto Político-Pedagógico os projetos integradores comuns a toda a escola, de acordo com o contexto de cada comunidade escolar, sendo estes planejados coletivamente pelos professores de cada \textbf{uma} das áreas do conhecimento.
Rearranjar os conhecimentos e as habilidades específicas de cada componente curricular e articulá -las aos conceitos estruturantes de \textbf{uma} área de conhecimento requer dois exercícios : o primeiro refere -se à necessidade de conhecimento aprofundado dos conteúdos, conceitos, proposições, modelos, esquemas, \textbf{teorias} e paradigmas \textbf{que} compõem o seu componente curricular e/ou especialidade.
O segundo exercício \textbf{diz} respeito à disposição para o diálogo, na busca de encontrar convergências e divergências, respeitando especificidades, não \textbf{hierarquizando} os conhecimentos escolares.
É a relação \textbf{que} se estabelece nesses encontros entre o específico aprofundado, próprio de cada componente curricular, e os conhecimentos em diálogo, quando mobilizados em \textbf{uma} proposta de trabalho por área do conhecimento, \textbf{que} os textos a seguir buscam \textbf{explicitar} a forma de organização curricular para a formação geral básica do Ensino Médio do Território Catarinense.
De acordo com a BNCC, a organização curricular do Novo Ensino Médio deve contemplar, também, a \textbf{abordagem} de “ temas \textbf{contemporâneos} \textbf{que} afetam a vida humana em escala local, regional e global, preferencialmente de forma transversal e integradora ” ( BRASIL, 2019, p. 19 ).
Para tanto, o trabalho com os chamados Temas \textbf{Contemporâneos} Transversais, retomados em documento complementar à BNCC, publicados pelo Ministério da Educação e Cultura, visam favorecer a integração dos componentes curriculares, estabelecendo “ sua conexão com situações vivenciadas pelos estudantes em suas realidades, contribuindo para trazer contexto e \textbf{contemporaneidade} aos objetos do conhecimento descritos na BNCC ” ( BRASIL, 2019, p. 5 ).
Os Temas \textbf{Contemporâneos} Transversais, abordados ao longo de todo o currículo, tanto na Formação Geral Básica, quanto na Parte Flexível do Currículo, figuram como questões profícuas e \textbf{centrais}, correspondendo a questões importantes, \textbf{urgentes} e presentes sob várias formas na vida cotidiana, oportunizando a elaboração \textbf{conceitual} complexa em torno de temáticas como Economia, Meio Ambiente, Saúde, Cidadania e Civismo, Interculturalidade, Ciência e Tecnologia, e possibilitando o diálogo fecundo dos conteúdos e conceitos dos componentes curriculares nas áreas do conhecimento.
Nesse segmento, os TCTs representam possibilidades fecundas para o \textbf{explicitar} das ligações entre os diferentes componentes curriculares, de forma integrada, bem como para estabelecer conexão com situações vivenciadas pelos \textbf{sujeitos} de aprendizagem, contemplando suas diversidades e contribuindo para trazer contexto e \textbf{contemporaneidade} aos objetos do conhecimento descritos no Currículo do Território Catarinense.
Além de garantir o \textbf{preconizado} pela legislação em torno de algumas temáticas, estes temas possuem grande potencialidade no \textbf{que} se refere ao desenvolvimento de \textbf{uma} educação voltada à cidadania como princípio \textbf{norteador} de aprendizagens, materializando, ainda, a transversalização como forma de organização das práticas escolares, conforme disposto nas DCNs : > A transversalidade é entendida como \textbf{uma} forma de organizar o trabalho \textbf{didático-pedagógico} em \textbf{que} temas, eixos temáticos são integrados às \textbf{disciplinas}, às áreas ditas convencionais de forma a estarem presentes em todas elas.
A transversalidade difere -se da \textbf{interdisciplinaridade} e complementam -se ; ambas rejeitam a concepção de conhecimento \textbf{que} toma a realidade como algo estável, pronto e acabado.
A primeira se refere à dimensão \textbf{didático-pedagógica} e a segunda, à \textbf{abordagem} epistemológica dos objetos de conhecimento.
A transversalidade orienta para a necessidade de se instituir, na prática educativa, \textbf{uma} analogia entre aprender conhecimentos teoricamente \textbf{sistematizados} ( aprender sobre a realidade ) e as questões da vida real ( aprender na realidade e da realidade ).
Dentro de \textbf{uma} compreensão interdisciplinar do conhecimento, a transversalidade tem significado, sendo \textbf{uma} proposta didática \textbf{que} possibilita o tratamento dos conhecimentos escolares de forma integrada.
Assim, nessa \textbf{abordagem}, a gestão do conhecimento parte do pressuposto de \textbf{que} os \textbf{sujeitos} são agentes da arte de problematizar e interrogar, e buscam procedimentos interdisciplinares capazes de acender a chama do diálogo entre diferentes \textbf{sujeitos}, ciências, saberes e temas ( BRASIL, 2013, p. 29 ). \textbf{Frisa} -se, \textbf{aqui}, de acordo com o \textbf{preconizado} pela BNCC e pelas DCN ( 2013 ), \textbf{que} os Temas \textbf{Contemporâneos} Transversais constituem \textbf{uma} referência nacional obrigatória, devendo \textbf{atravessar} todas as áreas do conhecimento.
Portanto, no presente documento curricular, estes temas são contemplados tanto na Formação Geral Básica, a partir de habilidades dos componentes curriculares, quanto nas Trilhas de Aprofundamento, posto \textbf{que} \textbf{uma} variedade destes temas foi elencada como temática \textbf{central} para esta importante parte dos Itinerários Formativos, assegurando o trato desses temas de forma contextualizada.
A figura 2 \textbf{explicita} de \textbf{que} forma se dá a organização destes temas.
Figura 2 - Temas \textbf{Contemporâneos} Transversais na BNCC Fonte : BRASIL.
Temas \textbf{Contemporâneos} Transversais na BNCC : Contexto Histórico e Pressupostos Pedagógicos ( 2019 ).
Parte -se da premissa de \textbf{que} os processos de \textbf{educar} e aprender devem pautar -se em situações de aprendizagem \textbf{que} associem de forma clara os conhecimentos escolares com a realidade vivenciada pelos \textbf{sujeitos} da escola, \textbf{uma} vez \textbf{que} essas vivências “ são fenômenos \textbf{que} envolvem todas as dimensões do ser humano ” e, quando não se alcança esta multidimensionalidade, produz -se “ alienação e perda do sentido social e individual ” ( BRASIL, 2019, p. 4 ).
Sendo assim, o trabalho pedagógico deve ser potencializado na medida em \textbf{que} enfatiza temáticas \textbf{que} possuem \textbf{uma} finalidade de crítica social.
Diante desta concepção, assume -se a progressão não seriada, compreendendo \textbf{que} o processo de elaboração \textbf{conceitual} com o qual se trabalha não é definido pelo ano letivo, mas, sim, pelos processos didáticos-pedagógicos construídos na sala de aula.
Portanto, compreende -se \textbf{que} o trabalho com os objetos de conhecimento, habilidades, valores e atitudes necessárias no ensino médio caminha na perspectiva ascendente, mas não necessariamente com \textbf{divisão} entre séries.
Por esse motivo, optou -se, neste material, por não fazer \textbf{divisões} dos objetos de conhecimento e habilidades entre as séries/anos do ensino médio.
Respeitando, assim, a autonomia, \textbf{que} é própria de cada projeto pedagógico escolar, e possibilitando ao corpo docente do ensino médio fazer escolhas sobre o aprofundamento de cada objeto de conhecimento.
Muito mais importante do \textbf{que} \textbf{dizer} \textbf{que}"x"ou"y"conteúdo deva ser desenvolvido nesta ou naquela série, está a compreensão de \textbf{que} o processo de elaboração \textbf{conceitual} demanda \textbf{uma} organização curricular \textbf{que} enfatize os processos de pensamento ao mesmo tempo em \textbf{que} produz, cria e aplica os conhecimentos.
Não se trata, portanto, de desordem ou apenas de evitar repetições \textbf{que} objetivam a \textbf{memorização} dos conteúdos.
Trata -se, sim, de defender \textbf{uma} organização curricular \textbf{que} promova a complexidade, o caráter dinâmico e não linear da realidade e \textbf{que} admita, desta forma, \textbf{que} o corpo docente promova a aprendizagem por diferentes e variados caminhos em torno dos problemas e preocupações sociais e pessoais de seus estudantes.
Assim, por se constituir como trabalho coletivo e reflexivo, produzido pelos educadores da Rede Estadual de Ensino de Santa Catarina, este documento foi construído por \textbf{uma} multiplicidade de vozes, \textbf{que} expressam a pluralidade marcada pela diversidade \textbf{que} o caracteriza. \textbf{Frisa} -se \textbf{que} algumas tensões foram importantes para orientar as reflexões e proposições de cada \textbf{uma} das quatro áreas do conhecimento \textbf{que} compõem a Formação Geral Básica.
Como em qualquer esforço coletivo, buscou -se representar alguns consensos ( e, sem dúvida, vários dissensos ) obtidos sem a pretensão de \textbf{que} estes representem a “ verdade ”, mas \textbf{uma} síntese necessária e provisória sobre estes novos arranjos curriculares denominados Áreas do Conhecimento.
Os textos \textbf{ora} apresentados representam a síntese produzida coletivamente por redatores, professores colaboradores/elaboradores, consultores e equipe SED, cujos nomes se encontram relacionados em cada \textbf{um} dos documentos.
Reconhecendo \textbf{que} estes textos são resultado de múltiplas interpretações e traduções dadas a partir dos contextos das práticas das escolas de ensino médio da rede estadual, ressalta -se \textbf{que}, por força do processo pelo qual foram construídos, há pontos relevantes nos pluralismos \textbf{teórico-metodológicos}, nas contradições e até mesmo nas incoerências.
No entanto, tais questões não retiram a legitimidade das proposições construídas a muitas mãos e em tempos políticos, econômicos e sociais tão conturbados, \textbf{demarcados}, ainda, por \textbf{uma} pandemia sem precedentes históricos.
Objetivou -se construir \textbf{um} material \textbf{que} permita o diálogo, a reflexão, a troca de ideias, imprimindo nele o contexto e o percurso de cada unidade escolar.
Por fim, cabe aos educadores assumirem a responsabilidade pela formação geral básica de milhares de estudantes, tendo a certeza e a convicção \textbf{que}, como afirmou Hannah Arendt : > A educação é a posição em \textbf{que} decidimos se amamos o mundo o \textbf{bastante} para assumir a responsabilidade por ele e, pela mesma razão, salvá -lo da ruína \textbf{que}, a não ser pela renovação, a não ser pela vinda do novo e dos jovens, seria inevitável.
E a educação é também quando decidimos se amamos nossos filhos o \textbf{bastante} para não os expulsar do nosso mundo e deixar \textbf{que} façam o \textbf{que} quiserem e \textbf{que} se virem sozinhos, nem para arrancar de suas mãos as mudanças de empreender algo novo, algo imprevisto por nós ( ARENDT, 1977, p. 196 ).

% matched lemmas: abordagem, aqui, atravessar, bastante, central, conceitual, configurar, contemporaneidade, contemporâneo, contextualização, contudo, demarcar, didático-pedagógico, disciplina, divisão, dizer, educar, explicitar, frisar, hierarquizar, implicar, interdisciplinaridade, memorização, norteador, operar, ora, perder, preconizar, que, santo, sistematizar, sujeito, teoria, teórico-metodológico, um, uma, urgente
\end{textsample}
